\chapter[Dissipative drag]{Irreversible Entropic Stress and Dissipative Drag}
\label{ch:dissipation}

\paragraph{Learning outcomes.}
After this chapter the reader should be able to:
\begin{enumerate}
  \item split \(T_{\mu\nu}^{\mathrm{ent}}\) into reversible and
        irreversible pieces and state the thermodynamic meaning
        of the split;
  \item reduce the irreversible piece to a two-body drag
        \(\mathbf{a}_{\mathrm{irrev}}\) consistent with
        Proposition~\ref{prop:allowed-drag};
  \item recover the working law
        \(\mathbf{F}_{\mathrm{diss}}=-\eta\mu\mathbf{v}/d^4\)
        as a laboratory, not as a multipole theorem;
  \item compare that channel with quadrupole gravitational-wave
        losses and say which process dominates where;
  \item enter Chapter~15 with an ordinary-differential system
        whose energy bookkeeping is already fixed.
\end{enumerate}

Chapter~13 used only the static extra density. Time-dependent
cores change their mutual information, and the first law of
Chapter~6 then requires an additional, velocity-odd stress.
Conservation still applies to the sum.

\section{Reversible versus irreversible}
\label{sec:ch14-split}

On a static spatial slice the mutual-information potential
\eqref{eq:Uent-def} is a function of the instantaneous regions
\(A\) and \(B\). Its metric variation is conservative: the
associated force derives from \(U_{\mathrm{ent}}\) and reverses
under \(\mathbf{v}\to-\mathbf{v}\).

If the regions change in time---because two cores inspiral,
because a star is sheared across a tidal radius, because a
condensate defect is stretched---then
\begin{equation}
  \label{eq:I-dot}
  \frac{\mathrm{d}}{\mathrm{d}t}I(A:B)
  \neq 0
\end{equation}
even after the conservative force has been subtracted. The work
associated with that change is the irreversible channel.

\begin{definition}[Irreversible entropic stress]
  \label{def:Tirrev}
  Decompose
  \begin{equation}
    \label{eq:split-rev-irrev}
    T_{\mu\nu}^{\mathrm{ent}}
    =T_{\mu\nu}^{\mathrm{ent,rev}}
     +T_{\mu\nu}^{\mathrm{ent,irrev}},
  \end{equation}
  where the first summand is the metric variation of
  \(U_{\mathrm{ent}}\) at frozen velocities and the second is
  whatever remainder is required for
  Theorem~\ref{thm:conservation} to hold on a time-dependent
  configuration. The remainder is odd under
  local time reversal of the relative motion when the
  entanglement production is itself odd.
\end{definition}

The split is thermodynamic rather than geometric:
\(T_{\mu\nu}^{\mathrm{ent,rev}}\) stores recoverable mutual
information; \(T_{\mu\nu}^{\mathrm{ent,irrev}}\) accounts for
entanglement that is produced or erased along the motion and is
not recovered by reversing the worldlines.

\section{Worldline reduction}
\label{sec:ch14-worldline}

Consider two compact condensed cores with worldlines
\(\gamma_A,\gamma_B\), rest masses \(m_A,m_B\), reduced mass
\(\mu=m_A m_B/(m_A+m_B)\), and separation four-vector
\(r^\mu=x_A^\mu-x_B^\mu\), \(d=\sqrt{r^\mu r_\mu}\) in a local
inertial frame. The entropic action along the worldlines is
\begin{equation}
  \label{eq:Sent-wl}
  S_{\mathrm{ent}}
  =-\int\mathrm{d}\tau\,U_{\mathrm{ent}}(d,\ell_A,\ell_B,\dot d),
\end{equation}
where \(\ell_A,\ell_B\) are proper core radii and dots are proper
time of the reduced motion. Varying with respect to \(x_A^\mu\)
produces the four-force~\eqref{eq:four-force}.

At leading post-Newtonian order, with \(\ell_A,\ell_B\) frozen,
\begin{equation}
  \label{eq:force-PN}
  \mathbf{F}_A
  =-\nabla_{\mathbf{x}_A}U_{\mathrm{ent}}(d)
   -\eta^{\mu}(d,\dot d)\,\mathbf{v}
   +O(c^{-2}).
\end{equation}
The first term is Chapter~13. The second term is the drag. The
coefficient \(\eta^{\mu}\) is not free: Proposition~\ref{prop:allowed-drag}
requires an entanglement energy \(E_{\mathrm{ent}}(d,\dot d)\) such
that
\begin{equation}
  \label{eq:balance-again}
  \frac{\mathrm{d}}{\mathrm{d}t}
  \bigl(E_{\mathrm{mech}}+E_{\mathrm{ent}}\bigr)
  =\mathbf{F}_{\mathrm{diss}}\cdot\mathbf{v}+O(c^{-2}),
\end{equation}
and the force on \(B\) is equal and opposite up to radiation
reaction of the tensor block.

A linear-response expansion of \(U_{\mathrm{ent}}\) about adiabatic
motion gives
\begin{equation}
  \label{eq:linear-response}
  \mathbf{F}_{\mathrm{diss}}
  =-\Gamma(d)\,\mathbf{v},
  \qquad
  \Gamma(d)>0
\end{equation}
whenever increasing \(d\) at finite speed decreases \(I(A:B)\)
relative to the adiabatic value. The sign is fixed by the
positivity of relative entropy: non-adiabatic separation produces
less mutual information than the static configuration at the same
\(d\), and the missing information is paid for by mechanical work.

\section{The working exponent}
\label{sec:ch14-exponent}

Section~\ref{sec:ch6-MI} gave the conformal large-separation tail
\begin{equation}
  \label{eq:I-tail}
  I(A:B)\;\sim\;c\,\frac{\ell_A^2\ell_B^2}{d^{2}}.
\end{equation}
Two derivatives with respect to time, or one derivative converting
a potential \(\sim 1/d^{2}\) into a velocity-dependent force after
linear response, produce a coefficient
\begin{equation}
  \label{eq:Gamma}
  \Gamma(d)
  =\eta\,\frac{\mu}{d^{4}}
\end{equation}
as a working choice: \(\eta\) absorbs \(c\), \(\ell_A,\ell_B\), and
the linear-response time. The source notes adopt this exponent.

\begin{assumption}[Working dissipative law]
  \label{ass:drag}
  In the reduced two-body laboratory of this book,
  \begin{equation}
    \label{eq:Fdiss}
    \mathbf{F}_{\mathrm{diss}}
    =-\eta\,\frac{\mu}{d^{4}}\,\mathbf{v},
  \end{equation}
  with \(\eta\ge 0\) constant on a given inspiral and with
  \(E_{\mathrm{ent}}\) defined so that~\eqref{eq:balance-again}
  holds identically.
\end{assumption}

Assumption~\ref{ass:drag} is not a multipole theorem. A different
tail of \(I(A:B)\) produces a different power. Chapter~16 treats
the exponent as falsifiable. Chapter~15 integrates~\eqref{eq:Fdiss}
and no other drag.

An explicit energy that satisfies the balance for
constant \(\eta\) is
\begin{equation}
  \label{eq:Eent-lab}
  E_{\mathrm{ent}}
  =-\frac{\eta\mu}{3d^{3}}
   \log\Bigl(1+\frac{|\mathbf{v}|^2}{v_*^2}\Bigr)
   +E_{\mathrm{ent}}^{\mathrm{(rev)}}(d),
\end{equation}
or any equivalent bookkeeping that makes
\(\dot E_{\mathrm{ent}}-\mathbf{F}_{\mathrm{diss}}\cdot\mathbf{v}\)
a total derivative at the post-Newtonian order kept. The integrator
in Chapter~15 may instead accumulate
\begin{equation}
  \label{eq:Wdiss}
  W_{\mathrm{diss}}(t)
  =\int_0^t\mathbf{F}_{\mathrm{diss}}\cdot\mathbf{v}\,dt'
\end{equation}
and display \(E_{\mathrm{mech}}(t)-W_{\mathrm{diss}}(t)\) as the
conserved mechanical-plus-entropic combination. Either
presentation is legal; omitting both is not.

\section{Comparison with quadrupole losses}
\label{sec:ch14-quad}

The Isaacson flux of the tensor block for a weak, slow binary is
the Peters--Mathews power
\begin{equation}
  \label{eq:Peters}
  P_{\mathrm{GW}}
  =\frac{32}{5}\Geff\mu^2 d^4\Omega^6
   \langle\text{inclination factors}\rangle
\end{equation}
in the usual quadrupole units, up to the precise harmonic
decomposition. The entropic power is
\begin{equation}
  \label{eq:Pent}
  P_{\mathrm{ent}}
  =-\mathbf{F}_{\mathrm{diss}}\cdot\mathbf{v}
  =\eta\mu\frac{v^2}{d^4}.
\end{equation}

For a circular Newtonian orbit, \(v^2=\Geff M/d\) and
\(\Omega^2=\Geff M/d^3\), so
\begin{equation}
  \label{eq:ratio}
  \frac{P_{\mathrm{ent}}}{P_{\mathrm{GW}}}
  \;\sim\;
  \frac{\eta}{\Geff^3 M^2\mu}\,
  \Bigl(\frac{d}{\Geff M}\Bigr)^{k}
\end{equation}
with a positive power of \(d\) that makes the ratio \emph{grow}
toward wide separations if \(\eta\) is held fixed. That conclusion
is an artefact of treating \(\eta\) as constant down to laboratory
distances. In a realistic condensate \(\eta\) itself falls with
the mutual-information tail and is negligible in the Solar System.
The intended domain of~\eqref{eq:Fdiss} is:

\begin{itemize}
  \item late inspirals of two compact condensed cores, where
        \(d\) is not astronomically large compared with
        \(\ell_A+\ell_B\);
  \item tidal disruption events, where a stellar core is sheared
        at \(d\sim r_{\mathrm{t}}\) and \(\dot d\) is not small;
  \item numerical laboratories that isolate the extra channel.
\end{itemize}

In LIGO-band vacuum inspirals of two objects whose condensate
cores are hidden behind horizons, quadrupole radiation remains
the default dissipative channel. NAMCG does not replace the
Peters formula by~\eqref{eq:Fdiss} in that regime unless
\(\eta\) is independently constrained to be large there. It is
not.

\section{Tidal disruption}
\label{sec:ch14-tde}

A star of radius \(R_*\) and mass \(m_*\) approaching a compact
object of mass \(M\) is disrupted at
\begin{equation}
  \label{eq:tidal-radius}
  r_{\mathrm{t}}
  \sim R_*\Bigl(\frac{M}{m_*}\Bigr)^{1/3}.
\end{equation}
Across that radius the baryonic body ceases to be a single region
\(A\) and becomes a collection of debris parcels. The mutual
information between the compact core and the star changes over a
dynamical time
\begin{equation}
  \label{eq:tdyn}
  t_{\mathrm{dyn}}
  \sim\sqrt{r_{\mathrm{t}}^3/\Geff M}.
\end{equation}
The irreversible channel is then unsuppressed: the process is
not adiabatic with respect to the modular Hamiltonian of the
combined system.

The book does not simulate debris streams. It records the
structural prediction that an extra, conserved drag of the type
\eqref{eq:Fdiss} can act near \(r_{\mathrm{t}}\) in addition to
ordinary hydrodynamics and tidal gravity. Observable consequences
(changes to fallback rates, extra heating) belong to Chapter~16
and require hydrodynamics that this text does not contain.

\section{Illegal and legal integrators}
\label{sec:ch14-legal}

An integrator that evolves
\begin{equation}
  \label{eq:illegal}
  \mu\ddot{\mathbf{r}}
  =-\frac{\Geff M\mu}{d^2}\hat{\mathbf{r}}
   -\gamma\mathbf{v}
\end{equation}
with constant \(\gamma\) and no \(E_{\mathrm{ent}}\) is not a
NAMCG calculation, even if its trajectories look like inspirals.
It dissipates energy into nowhere.

An integrator that evolves
\begin{equation}
  \label{eq:legal}
  \mu\ddot{\mathbf{r}}
  =-\nabla U_{\mathrm{N}}(d)
   -\nabla U_{\mathrm{ent}}^{\mathrm{rev}}(d)
   -\eta\frac{\mu}{d^4}\mathbf{v}
\end{equation}
and accumulates~\eqref{eq:Wdiss} is a legal laboratory of
Assumption~\ref{ass:drag}. Chapter~15 implements~\eqref{eq:legal}
with a symplectic Verlet step on the conservative forces and an
explicit treatment of the drag. The gravitational constant used
there is a toy value, not \(\Geff\) from Theorem~\ref{thm:Geff}.

\section{Exercises}
\label{sec:ch14-exercises}

\begin{enumerate}
  \item \textbf{Sign of \(\Gamma\).}
    From positivity of relative entropy, argue that
    non-adiabatic separation yields
    \(I(A:B)\le I_{\mathrm{adiabatic}}(d)\). Deduce
    \(\Gamma(d)\ge 0\) in~\eqref{eq:linear-response}.
  \item \textbf{Power counting.}
    If \(I\sim d^{-p}\), what exponent \(n\) appears in
    \(\Gamma(d)\propto d^{-n}\) after one linear-response time
    derivative? For which \(p\) does one recover \(n=4\)?
  \item \textbf{Energy identity.}
    For~\eqref{eq:Fdiss} with circular motion at instantaneous
    radius \(d(t)\), compute \(P_{\mathrm{ent}}\). Compare with
    \(\mathrm{d}U_{\mathrm{N}}/\mathrm{d}t\) along a Peters
    inspiral and estimate a separation where the two powers
    could be comparable for a chosen \(\eta\).
  \item \textbf{Equal and opposite.}
    Show that~\eqref{eq:Fdiss}, written in the reduced-mass
    frame, corresponds to equal-and-opposite three-forces on
    \(A\) and \(B\). What radiation-reaction term would have
    to be added to break that balance, and which block of
    \(T_{\mu\nu}^{\mathrm{total}}\) supplies it?
  \item \textbf{Constant-\(\gamma\) counterexample.}
    Take~\eqref{eq:illegal} and send \(d\to\infty\). What
    happens to the dissipated energy per unit time, and why
    is that incompatible with a mutual-information tail that
    vanishes at infinity?
  \item \textbf{Tidal time.}
    Evaluate~\eqref{eq:tdyn} for a solar-type star at the tidal
    radius of a \(10^6 M_\odot\) compact object. Why is the
    adiabatic approximation of Chapter~13 unreliable there?
\end{enumerate}

\section{Bridge to Chapter~15}
\label{sec:ch14-bridge}

The ordinary-differential laboratory is now specified:
Newtonian gravity, optional reversible extra potential, drag
\eqref{eq:Fdiss}, and an accumulated work~\eqref{eq:Wdiss}.
Chapter~15 integrates that system with a symplectic Verlet step,
plots the trajectory and the energy balance, and states, in the
first paragraph, that the figure is a reduced model. The Python
code that produces the figure lives in \texttt{code/} and is not
a derivation of Einstein gravity. After the figure, Chapter~16
collects every named assumption in the book and lists the
observations that would kill the framework.